\section{Reconstruction in the simplex Hilbert space}
\label{acfr:section}

The distribution hypothesis gives a uniform reconstruction of the coefficient
quotient in one dimension.
Divisor-density diagonalization gives the required norm estimate on every
coefficient direction.  Ordered intervals then place those directions in
$H_1=L^2([0,1],dt;\mathbb R)$.
Cell averages extend the reverse map to this Hilbert space.
In higher dimensions, controlled subspaces give a corresponding completion,
while concentrated samples obstruct a direct extension of point sampling.

\subsection{A uniform divisor-density estimate}

First take $k=1$ and one fixed nonnegative shift $h$.
Use the finite parameters of Theorem~\ref{acfd:exact}, and abbreviate
$\mathcal D=\{r<R:\mu^2(r)=1,\ (r,W)=1\}$.
For $y,z\in\mathbb R^{\mathcal D}$, put
\[
 D(y,z)=\sum_{r\in\mathcal D}\frac{y_rz_r}{\varphi(r)},\qquad
 \ell=1+\log R,\qquad d_N=\frac{eWR^2\ell^3}{N}.
\]
Here $d_N$ is the relative counting error for the one-dimensional reconstruction.

\begin{proposition}[Uniform mass comparison in one dimension]
For\label{acfr:density} all $y,z\in\mathbb R^{\mathcal D}$,
\begin{align}
 \sum_{d,e\in\mathcal D}\frac{(Ly)_d(Lz)_e}{[d,e]}
       &=D(y,z),\label{acfr:diagonal}\\
 \left|S_1(Ly,Lz)-\frac NW D(y,z)\right|
       &\leq eR^2\ell^3\sqrt{D(y,y)D(z,z)}.\label{acfr:relative-mass}
\end{align}
If $d_N<1$, the incidence map $T$ is injective.
For $W=W_N$ and $R=N^\rho$, with fixed $0<\rho<1/2$, one has $d_N\to0$.
\end{proposition}

\begin{proof}
The identity $\gcd(d,e)=\sum_{r\mid d,e}\varphi(r)$ gives
\[
 \sum_{d,e}\frac{\lambda_d\gamma_e}{[d,e]}
 =\sum_r\varphi(r)
       \left(\sum_{r\mid d}\frac{\lambda_d}{d}\right)
       \left(\sum_{r\mid e}\frac{\gamma_e}{e}\right).
\]
The inverse transform \eqref{eq:mobius-U} makes the right side $D(y,z)$
when $\lambda=Ly$ and $\gamma=Lz$.
All indices remain in the divisor-closed set $\mathcal D$.

Writing $\sigma_1(r)=\sum_{d\mid r}d$, the direct transform gives
\[
 \|Ly\|_1\leq\sum_{r\in\mathcal D}
       \frac{|y_r|\sigma_1(r)}{\varphi(r)}
 \leq\sqrt{D(y,y)}
       \left(\sum_{r<R}\frac{\sigma_1(r)^2}{\varphi(r)}\right)^{1/2}.
\]
The elementary bounds
$\sigma_1(r)\leq r(1+\log r)$ and
$r/\varphi(r)\leq e(1+\log r)$ imply
\[
 \sum_{r<R}\frac{\sigma_1(r)^2}{\varphi(r)}
       \leq e\ell^3\sum_{r<R}r\leq eR^2\ell^3.
\]
Apply \eqref{acfd:mass-spine} and \eqref{acfr:diagonal}.
If $TLy=0$ and $d_N<1$, then
$0\geq (N/W)(1-d_N)D(y,y)$, so $y=0$.
Invertibility of $L$ proves injectivity of $T$.
Finally $W_N=N^{o(1)}$ and $2\rho<1$ give $d_N\to0$.
\end{proof}

\subsection{The intervals for the divisor values}

The interval lengths require a uniform estimate for their total mass.
For squarefree $W\geq1$ and real $x\geq1$, set
\[
 G_W(x)=\sum_{\substack{r\leq x\\(r,W)=1}}
               \frac{\mu^2(r)}{\varphi(r)},\qquad
 E(W)=e^2\frac W{\varphi(W)}(10+\log W).
\]

\begin{lemma}[A uniform weighted harmonic sum]\label{acfr:harmonic}
For every such $W,x$,
\begin{equation}\label{acfr:harmonic-bound}
 \left|G_W(x)-\frac{\varphi(W)}W\log x\right|\leq E(W).
\end{equation}
\end{lemma}

\begin{proof}
Let $b_W$ be multiplicative, with $b_W(1)=1$.
For $p\nmid W$, its nonzero prime-power values are
\[
 b_W(p)=\frac1{p(p-1)},\qquad b_W(p^2)=-\frac1{p(p-1)}.
\]
For $p\mid W$, its only nonzero positive prime-power value is
$b_W(p)=-1/p$.
Checking prime powers proves the convolution identity
\[
 \frac{\mu^2(n)\mathbf1_{(n,W)=1}}{\varphi(n)}
       =\sum_{d\mid n}\frac{b_W(d)}{n/d}.
\]
Absolute Euler products give
\[
 A_W:=\sum_d|b_W(d)|
 \leq e^2\prod_{p\mid W}(1+1/p)
 \leq e^2\frac W{\varphi(W)},\qquad
 \sum_d b_W(d)=\frac{\varphi(W)}W.
\]
For the logarithmic moment, expand the nonnegative Euler products and
sum one prime contribution at a time.  This gives
\[
 \sum_d|b_W(d)|\log d
 \leq A_W\left(\sum_{p\mid W}\frac{\log p}{p}
          +3\sum_p\frac{\log p}{p(p-1)}\right)
 \leq A_W(\log W+9).
\]
The last prime sum is smaller than
$2\sum_{n\geq2}(\log n)/n^2<3$.
The final inequality follows by comparing this decreasing summand with
its integral from $2$ to infinity.
Truncation and monotone convergence justify the expansions.

Write $H_m=\sum_{j=1}^m1/j$.
The convolution identity yields
$G_W(x)=\sum_{d\leq x}b_W(d)H_{\lfloor x/d\rfloor}$.
For $u\geq1$, one has $0\leq H_{\lfloor u\rfloor}-\log u\leq1$.
Subtract $(\sum_db_W(d))\log x$.
The terms with $d\leq x$ are bounded by $|b_W(d)|(1+\log d)$.
For $d>x$, use $\log x\leq\log d$.
The preceding two moment bounds prove \eqref{acfr:harmonic-bound}.
\end{proof}

Return to the one-dimensional asymptotic parameters and write
\[
 B_N=\frac{\varphi(W_N)}{W_N}\log R_N,\qquad
 G_N=\sum_{r\in\mathcal D_N}\frac1{\varphi(r)},\qquad
 C_N=\frac N{W_N}B_N.
\]
Order $\mathcal D_N$ increasingly and put $x_r=\log r/\log R_N$.
Partition $[0,1)$ into consecutive half-open intervals $J_{r,N}$ of lengths
\[
 q_{r,N}=\frac1{\varphi(r)G_N}.
\]
Their endpoints follow the increasing order of $r$.
Define $D_Ny\in H_1$ to have value $y_r$ on $J_{r,N}$.
The symbol $D_N$ here denotes this step-function map.

\begin{lemma}[Location of the intervals]\label{acfr:cells}
One has $G_N/B_N\to1$.
With
\[
 h_N=\frac{2E(W_N)+1}{G_N},
\]
every $t\in J_{r,N}$ satisfies $|t-x_r|\leq h_N$.
In particular, $h_N\to0$ and the largest interval length tends to zero.
\end{lemma}

\begin{proof}
The strict cutoff $r<R_N$ changes $G_{W_N}(R_N)$ by at most one.
Thus $|G_N-B_N|\leq E(W_N)+1$.
The bound $W/\varphi(W)\leq e(1+\log W)$ gives
\[
 \frac{E(W_N)+1}{B_N}
 \ll\frac{(1+\log W_N)^3}{\log R_N}\longrightarrow0.
\]
Indeed, $\log W_N\leq D_0\log D_0$ and $\log R_N=\rho\log N$.

Let $H_N(t)=\sum_{x_r\leq t}q_{r,N}$ for $0\leq t\leq1$.
For $t<1$, its numerator is $G_{W_N}(R_N^t)$.
Equation~\eqref{acfr:harmonic-bound} therefore gives
$|H_N(t)-t|\leq h_N$.
The same bound holds at $t=1$, since $H_N(1)=1$.
The left and right endpoints of $J_{r,N}$ are $H_N(x_r^-)$ and $H_N(x_r)$.
For $r=1$, the left endpoint is zero.
For every other interval, take a left limit to obtain the claimed location.
\end{proof}

\subsection{The full one-dimensional reconstruction}

Let $\Pi_N=\#\{n\in I_N:n+h\in\mathbb P\}$ and put
\[
 \kappa_N=\frac{\Pi_NW_NG_N}{\rho N},\qquad
 \beta_N=\left(\frac{\kappa_NG_N}{B_N}\right)^{1/2}.
\]
The following statements use the distribution hypothesis \eqref{eq:lod}
and the fixed parameters of Theorem~\ref{thm:forms}.

\begin{theorem}[Reconstruction of the full quotient in one dimension]
\label{acfr:full-one}
For all sufficiently large $N$, the map
\begin{equation}\label{acfr:reconstruction}
 R_N^\sharp:\mathcal Q_N\longrightarrow H_1,\qquad
 R_N^\sharp[\lambda]=\beta_N D_N(U_N\lambda)
\end{equation}
is defined and injective.  Moreover, $\kappa_N\to1$, $\beta_N\to1$, and
\begin{align}
 \frac{\kappa_N}{1+d_N}\|z\|_N^2
 &\leq\|R_N^\sharp z\|_2^2
 \leq\frac{\kappa_N}{1-d_N}\|z\|_N^2,
                                      \label{acfr:norm-control}\\
 \frac{S_2(\lambda,\gamma)}{\rho C_N}
 &=\langle R_N^\sharp[\lambda],K_1R_N^\sharp[\gamma]\rangle.
                                      \label{acfr:full-prime}
\end{align}
For each fixed $F\in\mathcal X_1$,
$R_N^\sharp[E_NF]\to F$ in $H_1$.
The unrestricted normalized maximum satisfies
\begin{equation}\label{acfr:full-maximum}
 \frac{\kappa_N}{1+d_N}\leq\mathcal R_N
       \leq\frac{\kappa_N}{1-d_N},\qquad
 \mathcal R_N\longrightarrow1=M_1.
\end{equation}
\end{theorem}

\begin{proof}
Proposition~\ref{acfr:density} makes $T_N$ injective eventually.
Thus the quotient has a unique coefficient representative at those indices.
For $y=U_N\lambda$,
\[
 \|D_Ny\|_2^2=\frac{D(y,y)}{G_N},\qquad
 \int_0^1D_Ny=\frac{\lambda_1}{G_N}.
\]
The second identity follows from $(L_Ny)_1=\sum_ry_r/\varphi(r)$.
Every prime translate $n+h$, with $n\in I_N$, exceeds every supported divisor other than one.
Consequently $S_2(\lambda,\gamma)=\Pi_N\lambda_1\gamma_1$.
These identities prove \eqref{acfr:full-prime} whenever $\Pi_N>0$.
The mass comparison proves \eqref{acfr:norm-control}.

It remains to determine the scalar $\kappa_N$.
Choose a fixed nonnegative $F\in\mathcal X_1$ with $\int F>0$.
The interval-location lemma implies
$D_NP_NF\to F$ uniformly, since $F$ is uniformly continuous.
Hence
\[
 \frac{\bigl((E_NF)_1\bigr)^2}{B_NG_N}
       =\frac{G_N}{B_N}\left(\int_0^1D_NP_NF\right)^2
       \longrightarrow\left(\int_0^1F\right)^2>0.
\]
The exact prime identity writes $b_N(F,F)$ as $\kappa_N$ times this expression.
Theorem~\ref{thm:forms} gives the same positive limit for $b_N(F,F)$.
Division proves $\kappa_N\to1$, and also $\Pi_N>0$ eventually.
The location lemma now gives $\beta_N\to1$ and
$\beta_ND_NP_NF\to F$ for every fixed $F\in\mathcal X_1$.

For arbitrary $y$, Cauchy--Schwarz gives
$\lambda_1^2\leq G_ND(y,y)$.
Combine this with the lower mass bound to obtain the upper bound in
\eqref{acfr:full-maximum}.
The vector $y_r=1$ attains equality in that Cauchy--Schwarz bound.
The upper mass estimate applied to this vector gives the lower bound.
\end{proof}

For $k=1$, Proposition~\ref{prop:full-criterion} therefore applies with
\[
 u_N=\max\{0,\kappa_N/(1-d_N)-1\},\qquad v_N=0.
\]

\subsection{Cell averages and the Hilbert completion}

Let $\mathsf P_N$ be the orthogonal projection in $H_1$ onto the step functions
on the intervals $J_{r,N}$.
For $f\in H_1$, define
\begin{equation}\label{acfr:averages}
 S_N^\sharp f=[L_Ny(f)],\qquad
 y_r(f)=\frac1{\beta_Nq_{r,N}}\int_{J_{r,N}}f(t)\,dt.
\end{equation}
The integrals depend only on the $L^2$ class.

\begin{corollary}[Compatible maps on the completion]\label{acfr:completion}
For all sufficiently large $N$,
\[
 R_N^\sharp S_N^\sharp=\mathsf P_N,\qquad
 S_N^\sharp R_N^\sharp=I_{\mathcal Q_N},\qquad
 \|S_N^\sharp\|^2\leq\frac{1+d_N}{\kappa_N}.
\]
The projections $\mathsf P_N$ converge strongly to $I_{H_1}$.
For every fixed $F\in\mathcal X_1$,
\[
 \|S_N^\sharp F-[E_NF]\|_N\longrightarrow0.
\]
For arbitrary $f,g\in H_1$, the exact prime-boundary pairing is
\begin{equation}\label{acfr:completion-pairing}
 \widehat\mu_1\bigl(\Theta_N(S_N^\sharp f,S_N^\sharp g;1)\bigr)
       =\langle f,K_1g\rangle.
\end{equation}
Here $\Theta_N$ acts on quotient classes as in Corollary~\ref{acfp:quotient}.
\end{corollary}

\begin{proof}
Cell averaging proves both composition identities.
Apply \eqref{acfr:norm-control} to $S_N^\sharp f$ and use
$\|\mathsf P_Nf\|_2\leq\|f\|_2$ to obtain the norm bound.
The largest interval length tends to zero.
Uniform continuity first gives $\mathsf P_Nf\to f$ for continuous $f$.
The contraction bound and density of continuous functions extend this to $H_1$.

For fixed smooth $F$, both
$R_N^\sharp[E_NF]$ and $R_N^\sharp S_N^\sharp F$ tend to $F$.
The lower bound in \eqref{acfr:norm-control} proves their coefficient-norm comparison.
Finally, $\mathsf P_N$ preserves the integral, and
$K_1g=(\int_0^1g)1$.
Combine \eqref{acfr:full-prime} with the exact boundary pairing to obtain
\eqref{acfr:completion-pairing}.
\end{proof}

This also specifies which coefficient sequences define a vector in the completion.
Take sequences $z_N\in\mathcal Q_N$ for which $R_N^\sharp z_N$ is Cauchy in $H_1$.
Identify two such sequences when $\|z_N-z'_N\|_N\to0$.
The map
\[
 [(z_N)]\longmapsto\lim_NR_N^\sharp z_N
\]
is an isometric bijection onto $H_1$, with inner product
$\lim_N\langle z_N,z'_N\rangle_N$.
Indeed, \eqref{acfr:norm-control} and polarization identify that limit.
Every $f\in H_1$ has the representative $z_N=S_N^\sharp f$.
Boundedness of $(z_N)$ alone does not imply the required Cauchy property.

For a real counting observable $w:I_N\to\mathbb R$, put
\[
 m_N(w)=\frac1{\Pi_N}\sum_{\substack{n\in I_N\\n+h\in\mathbb P}}w(n).
\]
Since each counted prime translate has incidence value $\lambda_1$, the same proof gives
\begin{equation}\label{acfr:weighted-one}
 \widehat\mu_1\bigl(\Theta_N(S_N^\sharp f,S_N^\sharp g;w)\bigr)
       =m_N(w)\langle f,K_1g\rangle,
 \qquad |m_N(w)|\leq\|w\|_\infty.
\end{equation}
For a uniformly bounded sequence of observables, these forms converge on all
of $H_1$ exactly when the scalar sequence $m_N(w_N)$ converges.
In that case the convergence is in operator norm.
Testing $f=g=1$ proves necessity.

\subsection{Two independent coefficient directions in dimension two}

For $k=2$, take $\mathcal H=\{0,2\}$, $N=1024$, $R=16$, $W=2$, and $v=1$.
The only sampling nodes in the interior of $\Delta_2$ are
\[
 x=\frac{(\log3,\log5)}{\log16},\qquad
 x'=\frac{(\log5,\log3)}{\log16}.
\]
Both have coordinate product $15<R$ before the logarithmic change.
Disjoint interior smooth bumps give sample vectors supported on either node.
Let $\lambda=L\mathbf1_{(3,5)}$ and $\gamma=L\mathbf1_{(5,3)}$.
Their incidence functions are
\[
 u(n)=\frac{(1-3\mathbf1_{3\mid n})(1-5\mathbf1_{5\mid n+2})}{8},\qquad
 v(n)=\frac{(1-5\mathbf1_{5\mid n})(1-3\mathbf1_{3\mid n+2})}{8}.
\]
For $w_0(n)=1$ and the independent counting observable $w_1(n)=n/N$, define
\[
 B_j=\sum_{n\in I_N}w_j(n)R_{\mathcal H}(n)
       \begin{pmatrix}u(n)^2&u(n)v(n)\\u(n)v(n)&v(n)^2\end{pmatrix}.
\]
Direct finite summation gives
\begin{equation}\label{acfr:two-matrices}
 B_0=\frac1{64}\begin{pmatrix}1018&73\\73&976\end{pmatrix},\qquad
 B_1=\frac1{65536}\begin{pmatrix}1542766&111387\\111387&1508932\end{pmatrix}.
\end{equation}
The corresponding sums without $R_{\mathcal H}$ are
\[
 A_0=\frac1{32}\begin{pmatrix}2041&252\\252&2050\end{pmatrix},\qquad
 A_1=\frac1{32768}\begin{pmatrix}3145686&387073\\387073&3151344\end{pmatrix}.
\]
These values follow by testing the odd integers from $1025$ to $2047$.
Trial division through $43$ determines primality of both translates,
since $47^2>2049$.
Both prime matrices have rank two.  In fact,
\[
 \det B_0=\frac{988239}{4096}>0,\qquad
 \det B_1=\frac{2315521922143}{4294967296}>0.
\]
The matrices are not proportional, since
\[
 (B_1)_{11}(B_0)_{12}-(B_1)_{12}(B_0)_{11}=-\frac{47}{256}.
\]
Here $\rho=2/5$, $C_N=2048(\log2)^2$, and the boundary form matrix is
$5B_j/[4096(\log2)^2]$.
Thus a scalar prime mean multiplying one rank-one form cannot represent
these two coefficient directions and their counting observable.

\subsection{A Cauchy-sequence obstruction for raw sampling}

The reconstruction in dimension one does not extend the point-sampling map
to arbitrary $L^2$ equivalence classes.
The following two-dimensional sequence makes that distinction quantitative.
Choose $b\in C_c^\infty((1/8,3/8))$ with $0\leq b\leq1$ and
$b=1$ on $[3/16,5/16]$.
For the cutoff in Lemma~\ref{acfd:cutoff}, one such choice is
$b(t)=\chi(32t-4)\chi(12-32t)$.
Set $F_0(t_1,t_2)=b(t_1)b(t_2)$.
Its support lies in the interior of $\Delta_2$, and
\begin{equation}\label{acfr:positive-model}
 \langle F_0,K_2F_0\rangle
       =2\left(\int b\right)^2\int b^2\geq\frac1{256}.
\end{equation}

\begin{proposition}[Vanishing Hilbert norm with nonvanishing sampled pairing]
\label{acfr:cauchy-obstruction}
Assume \eqref{eq:lod} for the two fixed shifts $\{0,2\}$ and the asymptotic
parameters of Theorem~\ref{thm:forms}.  There are explicit $F_N\in\mathcal X_2$ such that
\[
 \|F_N\|_{H_2}\leq N^{-1},\qquad E_NF_N=E_NF_0.
\]
For $w_N(n)=n/N$, their exact boundary pairings satisfy, eventually,
\begin{equation}\label{acfr:nonvanishing}
 \widehat\mu_1\bigl(\Theta_N(E_NF_N,E_NF_N;w_N)\bigr)
       \geq\frac1{512}.
\end{equation}
Thus raw point sampling followed by the normalized prime pairing has no
uniformly bounded extension to the $H_2$ unit balls.
The sequence $F_N$ tends to zero in the original Hilbert completion.
At each sufficiently large fixed $N$, the form also has no bounded extension to $H_2$.
\end{proposition}

\begin{proof}
Let $X_N$ be the finite set of distinct sampling nodes.
Let $X_N^+=\{x\in X_N:F_0(x)\ne0\}$ and $m_N=|X_N^+|$.
The product-count bound gives $m_N\leq R_N(1+\log R_N)$.
Distinct nodes differ in some coordinate by at least
$1/(R_N\log R_N)$, since $\log(b/a)\geq1/b$ for integers $1\leq a<b<R_N$.
Choose a fixed $\psi\in C_c^\infty((-1,1)^2)$ with
$0\leq\psi\leq1$ and $\psi(0)=1$.
For example, take $\psi(u_1,u_2)=q(u_1)q(u_2)$, where
$q(u)=\exp(1-(1-4u^2)^{-1})$ for $|u|<1/2$ and $q(u)=0$ otherwise.
Define
\[
 a_N=\min\left\{\frac1{32},\frac1{4R_N\log R_N},
              \frac1{2N\sqrt{R_N(1+\log R_N)}}\right\}>0,\qquad
 F_N(t)=\sum_{x\in X_N^+}F_0(x)\psi\left(\frac{t-x}{a_N}\right).
\]
The supports are disjoint and lie inside $\Delta_2$.
Each contains just its central sampling node.
Thus $P_NF_N=P_NF_0$, and the coefficient identity follows.
Each bump has area at most $4a_N^2$ and absolute value at most one.
Consequently $\|F_N\|_2^2\leq4m_Na_N^2\leq N^{-2}$.

For $n\in I_N$, one has $w_N(n)\geq1$.
Positivity of the squared incidence weights bounds the pairing below by
$S_2(E_NF_0)/(\rho C_N)$.
Theorem~\ref{thm:forms} makes this converge to
$\langle F_0,K_2F_0\rangle\geq1/256$.
This proves \eqref{acfr:nonvanishing}.
Any common Hilbert-norm bound for these quadratic forms would instead
bound their values by a constant times $N^{-2}$.
For fixed $N$, replace the bump width by $a_N/j$ and let $j\to\infty$.
The samples remain fixed while the Hilbert norm tends to zero.
\end{proof}

This obstruction concerns the raw sampling operation on $H_2$.
It does not exclude bounded maps from the finite coefficient quotients to $H_2$.
The functions $F_N$ have increasingly large derivative norms, so the
uniform estimate on $\mathcal X_2$ remains consistent with the example.

\subsection{Reconstruction on controlled coefficient subspaces}

There is a completion compatible with the prime boundary in every dimension.
Use the nested spaces $V_{j(N)}$ from Corollary~\ref{cor:diagonal} and write
$V_N=V_{j(N)}$, $P_N^V=P_{V_N}$, and $G_N^V=G_{N,V_N}$.
For sufficiently large $N$, define
\begin{gather*}
 Z_N:V_N\to\mathcal Q_N,\qquad
 Z_Nu=[E_N(G_N^V)^{-1/2}u],\\
 \mathscr S_N=Z_NP_N^V:H_k\to\mathcal Q_N,\qquad
 \mathscr R_N=Z_N^*:\mathcal Q_N\to V_N\subset H_k.
\end{gather*}
The adjoint uses the quotient inner product $S_1/C_N$.

\begin{theorem}[The comparison on the Hilbert completion]
\label{acfr:regular-completion}
The map $Z_N$ is an isometry.  Hence
\[
 \|\mathscr R_N\|=\|Z_N\|=1,\qquad
 \|\mathscr S_N\|\leq1,\qquad
 \mathscr R_N\mathscr S_N=P_N^V,\qquad
 \mathscr S_N\mathscr R_N=\operatorname{proj}_{Z_N(V_N)}.
\]
For every $f,g\in H_k$,
\begin{equation}\label{acfr:regular-boundary}
 \widehat\mu_1\bigl(\Theta_N(\mathscr S_Nf,\mathscr S_Ng;1)\bigr)
       \longrightarrow\langle f,K_kg\rangle.
\end{equation}
For every fixed $F\in\mathcal X_k$,
$\|\mathscr S_NF-[E_NF]\|_N\to0$.

More generally, suppose $z_N,z'_N\in\mathcal Q_N$ have bounded norms and
\[
 \operatorname{dist}(z_N,Z_N(V_N))\to0,\qquad
 \operatorname{dist}(z'_N,Z_N(V_N))\to0,\qquad
 \mathscr R_Nz_N\to f,\quad\mathscr R_Nz'_N\to g.
\]
Their unweighted prime-boundary pairing converges to $\langle f,K_kg\rangle$.
\end{theorem}

\begin{proof}
By the definition of the Gram operator,
$\|Z_Nu\|_N^2=\langle (G_N^V)^{-1/2}u,G_N^V(G_N^V)^{-1/2}u\rangle
=\|u\|_2^2$.
This proves the isometry and the composition identities.
Put $Q_N^V=(G_N^V)^{-1/2}H_{N,V_N}(G_N^V)^{-1/2}$.
The exact boundary pairing on $Z_N(V_N)$ is
\[
 \widehat\mu_1(\Theta_N(Z_Nu,Z_Nv;1))=\langle u,Q_N^Vv\rangle.
\]
Theorem~\ref{thm:galerkin} gives
$\|Q_N^V-P_N^VK_k|_{V_N}\|\to0$.
Since $P_N^V\to I$ strongly, this proves \eqref{acfr:regular-boundary}.

If $F$ lies in the union of the $V_j$, then eventually $P_N^VF=F$.
The bound $\|G_N^V-I\|\to0$ gives
$\|\mathscr S_NF-[E_NF]\|_N\to0$.
For an arbitrary fixed $F\in\mathcal X_k$, choose a fixed
$v\in\bigcup_jV_j$ close to $F$ in $H_k$.
Theorem~\ref{thm:forms} gives
$\|[E_N(F-v)]\|_N\to\|F-v\|_2$, and $\mathscr S_N$ is a contraction.
The triangle inequality gives a limiting upper bound $2\|F-v\|_2$.
Density proves the stated comparison.

On the full quotient, Cauchy--Schwarz and $0\leq R_{\mathcal H}\leq k$ give
\begin{equation}\label{acfr:full-bound}
 \left|\frac{S_2(\lambda,\gamma)}{\rho C_N}\right|
       \leq\frac k\rho\|[\lambda]\|_N\|[\gamma]\|_N.
\end{equation}
Replacing each bounded sequence by its projection onto $Z_N(V_N)$
therefore changes the pairing by a quantity tending to zero.
The projected vectors are $Z_N\mathscr R_Nz_N$ and
$Z_N\mathscr R_Nz'_N$.
The finite operator estimate and convergence of their reconstructed vectors
prove the final assertion.
\end{proof}

For any uniformly bounded real counting observables $w_N$, the same maps
give exact forms on $H_k$ by pullback through $\mathscr S_N$.
Their norms are at most $(k/\rho)\sup_N\|w_N\|_\infty$.
A limit for a nonconstant observable requires convergence of these
compressed weighted forms on a dense set.
The uniform bound makes this condition sufficient for a unique bounded
limit form on $H_k$.
That additional convergence is not supplied by the unweighted form theorem.
Equation~\eqref{acfr:two-matrices} gives their first two-dimensional arithmetic values.

For $k\geq2$, the full-quotient question remains the control of directions
outside $Z_N(V_N)$.  The norm approximation assumed in
Theorem~\ref{acfr:regular-completion} is a stated restriction on sequences.
Neither that theorem nor the sampling obstruction establishes the two
uniform inequalities in Proposition~\ref{prop:full-criterion} on all of $\mathcal Q_N$.
The explicit-formula comparison also retains its separate test-domain and
convergence requirements.
